\documentclass[10pt]{article}
\usepackage[margin=0.8in]{geometry}
\usepackage{booktabs}
\usepackage{longtable}
\usepackage{hyperref}
\usepackage{xcolor}

\title{CARLA Highway RL: Technical and Iteration Log}
\author{Living implementation record}
\date{\today}

\begin{document}
\maketitle

% APPEND-ONLY EXPERIMENT RECORD:
% Preserve prior run entries, failures, configuration changes, and resolutions.
% Correct an old entry with a new dated note rather than deleting history.

\section{Runtime and CARLA Installation Notes}
The implementation targets Python 3.12 and the packaged CARLA 0.9.16 UE4.26
release. The CARLA Python client and simulator server are separate processes.
External mode is the default; managed mode requires an explicit CARLA root and
records the repository-owned process group before it can be stopped. The
implementation task did not install or run a CARLA server.

\section{Repository Architecture}
The exact three Python source files are \texttt{carla\_env.py},
\texttt{policies.py}, and \texttt{run.py}. Environment, controller, sensors,
traffic, and geometry remain cohesive in the first file; four tactical policy
adapters remain in the second; orchestration, statistics, plotting, video, sync,
and report generation remain in the third. Configuration is centralized in
\texttt{config.yaml}.

\section{Environment Lifecycle}
The environment connects, validates exact client/server versions, loads Town04
when required, saves original world settings, applies synchronous fixed-step
settings, and makes itself the only ticking client. Reset destroys prior
sensors before vehicles, increments an episode token to invalidate callbacks,
reseeds every random source, and rebuilds all actor references. Close is
idempotent and restores owned world settings.

\section{Highway Spawn Discovery}
Each map spawn is projected to a driving waypoint. Junctions, single-lane
locations, opposite-direction adjacency, and insufficient non-junction forward
continuity are rejected. Remaining candidates are deterministically scored by
adjacent-lane count, continuity, nearby spawn availability, and spawn index.
Doctor writes the ordered candidate records for inspection.

\section{Neighbor Classification}
Nearby vehicles are first matched by road and lane identity. A geometric
fallback uses longitudinal and lateral projection in the candidate lane frame
with a same-direction tangent check, covering road or section transitions.
Positive longitudinal projection denotes a front vehicle; deterministic actor
ordering resolves ties.

\section{Lane-Change Controller}
The tactical validator checks adjacent driving-lane legality, active
lane-change state, front and rear gaps, and rear TTC. Accepted changes retain a
target road/lane key and require lane identity plus small center offset for
several consecutive ticks before one completion is counted. Repeated tactical
requests cannot increment completion counts.

\section{Reward Iterations}
The implementation uses only progress, capped speed, rejected-request cost,
accepted-change cost, and terminal collision/off-road/stuck penalties. No
reward iteration or tuning experiment has been performed.

\section{PPO Runs}
No PPO runs were performed during repository implementation. Fresh and
additional-budget resume paths share Stable-Baselines3 PPO parameters from the
resolved configuration; resume retains timestep numbering.

\section{Baseline Runs}
No server-backed baseline episodes were performed during repository
implementation. Offline synthetic checks cover action validity and the blocked
lane branch of the rule-based policy.

\section{Bugs and Resolutions}
No experimental bugs are recorded. Implementation-time validation findings
should be appended here with a date once offline and CARLA-backed checks run.

\section{Evaluation Manifest Versions}
No manifest has been frozen yet. Each future manifest record must include its
SHA-256 hash, scenario count, seed set, and deliberate replacement reason, if
any.

\section{Video Rendering Notes}
The RGB callback uses a bounded queue and exact simulator-frame matching.
CARLA's BGRA buffer is converted to RGB for environment output and then to BGR
for OpenCV encoding. Rendering requires
\texttt{world.no\_rendering\_mode=false}.

\section{Reproducibility Checklist}
\begin{itemize}
  \item Resolved configuration and hash saved for training and evaluation.
  \item Python, NumPy, Traffic Manager, environment, and SB3 seeds fixed.
  \item Synchronous mode and 0.05\,s fixed delta applied.
  \item Evaluation conditions paired across all four policies.
  \item Client/server version mismatch fails before an episode.
  \item Generated results remain absent or TBD until experiments run.
\end{itemize}

\section{Optional Future Work}
The only documented optional extension is a separately labeled CARLA
BehaviorAgent comparison after the core four-policy experiment is complete.

\section{Experiment Ledger}
Append one row per material run or decision. Do not rewrite historical rows.

\begin{longtable}{p{1.2cm}p{1.6cm}p{1.2cm}p{0.7cm}p{1.1cm}p{1.4cm}p{1.3cm}p{1.6cm}p{1.5cm}}
\toprule
Date & Run name & Git hash & Seed & Timesteps & Config hash & Result & Issue & Decision \\
\midrule
\endhead
Implementation & repository build & pending & 0 & --- & generated at run time &
Source complete & CARLA unavailable in build environment &
Run offline checks, then doctor and smoke before training \\
\bottomrule
\end{longtable}

\end{document}
